\chapter{总结与展望}\label{chap:conclusion}

\section{全文总结}

本报告围绕八皇后问题完成了六类实现的统一对照：命令式穷举+逻辑筛选、纯逻辑 \texttt{permuteq}、自定义不等关系逻辑求解、DFS、BFS 与位运算回溯。主要结论如下。

\begin{enumerate}
    \item 六种方法都稳定得到 92 个合法解，正确性一致；
    \item 纯逻辑 \texttt{permuteq} 在当前环境下耗时最高；
    \item 在逻辑路线中加入自定义不等关系并前置约束后，性能有明显改善；
    \item 位运算回溯耗时最低，约为 DFS 的 $9.76\times$ 速度、BFS 的 $9.59\times$ 速度。
\end{enumerate}

如果把这些结果压缩成一句话，就是：\textbf{剪枝策略决定搜索规模，状态表示决定单节点开销}。这两个因素叠加后，才形成最终运行时间。

\section{展望}

下一步可以从三个方向继续推进：

\begin{enumerate}
    \item 扩展到更大规模 $N$，观察不同方法在规模增长下的拐点；
    \item 引入更强的约束传播和变量选择启发式，进一步压缩无效搜索；
    \item 在保持逻辑建模可解释性的前提下，探索更高性能的执行后端。
\end{enumerate}
